% 第5章第7节：本章小结。

\section{本章小结}
\label{sec:graph-dp-summary}

本章处理的问题由局部函数组成：每个变量取有限个状态，整体成本是若干局部项之和，每一项只读取少数变量。计算的核心始终是同一个动作：固定一组接口变量，把只在一侧出现的变量汇总掉，留下一张以接口为索引的表。链上，接口是一个变量，表沿链逐个传递；树上，接口仍是一个变量，但一个顶点要合并多棵子树传来的消息；一般依赖图上，变量消元按任意次序逐个汇总变量，正确性不变，中间表的作用域却随次序变化。树分解把这些同时处理的变量组织成袋树，消元次序可以机械地产生树分解，弦图上还能用极大团和最大权生成树得到团树；最后，袋上的消息把树上的递推推广到一般图。

这些步骤中有三种不能互相替代的对象。局部函数表记录数值，决定消息的具体取值；原始图记录哪些变量共同出现，决定消元时的填充边和袋必须容纳的团；袋树记录计算的组织方式，决定消息沿什么路径传递、以什么为索引。原始图不能恢复函数表，袋树的边也不是原始图的边；混淆它们会导致把高阶约束错拆成成对规则，或在袋之间重复累计成本。

精确计算依赖一组明确的条件：局部函数给出完整的作用域；状态集合有限；组合与汇总的运算满足结合律、交换律与分配律，即构成半环；每个局部函数恰好归入一个包含其作用域的袋；袋树满足运行交集，使接口把两侧的变量隔开。计算规模则主要由袋的大小控制，宽度为$k$、每个变量至多$q$个状态时，袋表至多$q^{k+1}$项。树宽小的图上精确计算可行，树宽大的图上同样的方法仍然正确，但可能无法承受。后续的概率模型会把这些表解释为未归一化的概率、把求和解释为边缘化，但归一化条件、查询中不同汇总的次序以及带环图上的近似方法，都需要在那里另外说明，不能由本章的结论直接推出。

本章比较的是有限配置：每个变量取离散的状态，目标是在所有配置中找最小值或求汇总。下一章转向另一类任务：顶点上的数值已经给定，要研究它们沿边怎样变化，怎样通过图上的线性算子对它们进行变换和滤波，并据此划分图。所需的工具也随之改变，从有限比较与递推转向内积、二次型和第三部分的谱分解。
